\section{Conjugation, the differential of the exponential, and Bernoulli series}\label{sec:differential}
This section opens a second route to the BCH series, independent of the Hopf
algebra argument. It also proves the Campbell identity and every
infinitesimal identity used later. Parameter differentiation of exponentials
is treated systematically in \cite{wilcox}; all formulas required here are
derived in full.

\subsection{Campbell's conjugation identity and the braiding formulas}
\begin{theorem}[Campbell--Hadamard identity]\label{thm:campbell}
Let either $X,Y$ have positive degree in $\widehat T$ with $s$ a formal or
scalar parameter, or let $X,Y$ lie in a Banach algebra with $s$ an
arbitrary scalar. Then
\begin{equation}\label{eq:campbell}
 e^{sX}Ye^{-sX}=e^{s\ad_X}Y=\sum_{n\geq0}\frac{s^n}{n!}\ad_X^nY.
\end{equation}
For invertible elements write $\Ad_AY=AYA^{-1}$; then
$\Ad_{e^X}=e^{\ad_X}$. The Banach-algebra series converges absolutely for
every $X,Y,s$, with
$\sum_n|s|^n\norm{\ad_X^nY}/n!\leq e^{2|s|\norm X}\norm Y$.
An entirely associative version is
\begin{equation}\label{eq:adbinomial}
 \ad_X^nY=\sum_{j=0}^{n}(-1)^j\binom nj X^{n-j}YX^j.
\end{equation}
\end{theorem}
\begin{proof}
Let $L_X$ and $R_X$ denote left and right multiplication by $X$. They
commute, since $L_XR_XA=XAX=R_XL_XA$, and $\ad_X=L_X-R_X$. The binomial
theorem for commuting operators gives
$\ad_X^n=\sum_j(-1)^j\binom njL_X^{n-j}R_X^j$, which applied to $Y$ is
\eqref{eq:adbinomial}.

Set $F(s)=e^{sX}Ye^{-sX}$. In a Banach algebra $F$ is differentiable with
$F'(s)=Xe^{sX}Ye^{-sX}-e^{sX}Ye^{-sX}X=\ad_XF(s)$ and $F(0)=Y$, by the
product rule and $\frac{\dd}{\dd s}e^{sX}=Xe^{sX}=e^{sX}X$. The map
$\ad_X$ is a bounded linear operator with $\norm{\ad_X}\leq2\norm X$, so the
linear differential equation $F'=\ad_XF$ has the unique solution
$F(s)=e^{s\ad_X}Y=\sum_ns^n\ad_X^nY/n!$; uniqueness holds because a
difference $G$ of two solutions satisfies $G(0)=0$ and
$\norm{G(s)}\leq2\norm X\bigl|\int_0^s\norm{G(\sigma)}\,\dd\sigma\bigr|$ on
every compact interval, hence vanishes by Gr\"onwall's inequality. The norm bound follows from $\norm{\ad_X^nY}\leq(2\norm X)^n\norm Y$.
Formally, both sides of \eqref{eq:campbell} are power series in $s$ with
coefficients in $\widehat T$; using \eqref{eq:adbinomial}, the coefficient
of $s^n$ on the left is
$\sum_{j}X^{n-j}Y(-X)^j/((n-j)!j!)=\frac1{n!}\sum_j\binom nj(-1)^jX^{n-j}YX^j
=\ad_X^nY/n!$, which is the coefficient on the right. (This direct comparison
also gives a second, purely algebraic proof in the Banach case.) Setting
$s=1$ gives $\Ad_{e^X}=e^{\ad_X}$.
\end{proof}

Conjugation $A\mapsto e^XAe^{-X}$ is an algebra automorphism; it is
continuous in the Banach case, and in the formal case it preserves the
filtration ($e^XF^m\widehat Te^{-X}\subseteq F^m\widehat T$, since $e^{\pm X}$
have constant term $1$) and is therefore continuous for the degree
topology. Hence it commutes with the exponential series. Therefore
\begin{align}
 e^Xe^Ye^{-X}&=\exp\bigl(e^XYe^{-X}\bigr)=\exp\bigl(e^{\ad_X}Y\bigr),\label{eq:conjexp}\\
 e^Xe^Y&=\exp\Bigl(\sum_{n\geq0}\frac{\ad_X^nY}{n!}\Bigr)e^X.\label{eq:braiding}
\end{align}
The second follows from the first by multiplying on the right by $e^X$. These
are the general braiding identities of the source page, including the
iterated-commutator notation $[X,Y]_n=\ad_X^nY$ there. Unlike a general BCH
series, the adjoint exponential in \eqref{eq:braiding} is an entire series
for bounded operators and requires no smallness.

\subsection{The differential of the exponential}
\begin{theorem}[Duhamel differential]\label{thm:dexp}
Let $A(t)$ be a differentiable curve in a Banach algebra, with $H=A'(t)$,
or a formal power series in $t$ with coefficients in $F^1\widehat T$. Then
\begin{align}
 (\dd\exp)_A(H):=\frac{\dd}{\dd t}e^{A(t)}
 &=\sum_{p,q\geq0}\frac{A^pHA^q}{(p+q+1)!}
 =\int_0^1e^{(1-s)A}He^{sA}\,\dd s,\label{eq:duhamel}\\
 e^{-A}\frac{\dd}{\dd t}e^{A}
 &=\int_0^1e^{-s\ad_A}H\,\dd s
 =\sum_{n\geq0}\frac{(-1)^n}{(n+1)!}\ad_A^nH=\phi(\ad_A)H,\label{eq:leftdexp}\\
 \Bigl(\frac{\dd}{\dd t}e^{A}\Bigr)e^{-A}
 &=\int_0^1e^{s\ad_A}H\,\dd s
 =\sum_{n\geq0}\frac{1}{(n+1)!}\ad_A^nH=\phi(-\ad_A)H.\label{eq:rightdexp}
\end{align}
All series converge absolutely for arbitrary bounded $A,H$.
\end{theorem}
\begin{proof}
The product rule, applied without moving any factor, gives
$\frac{\dd}{\dd t}A^n=\sum_{j=0}^{n-1}A^jHA^{n-1-j}$; in the Banach case the
exponential series may be differentiated term by term because
$\sum_n\norm{\frac{\dd}{\dd t}A^n}/n!\leq\sum_nn\norm A^{n-1}\norm H/n!$
converges uniformly on compact parameter sets, and in the formal case the
statement is coefficientwise. Collecting terms, the coefficient of $A^pHA^q$
in $\frac{\dd}{\dd t}e^A$ is $1/(p+q+1)!$, which is the first equality in
\eqref{eq:duhamel}. Expanding the two exponentials in the integral, the
coefficient of $A^pHA^q$ is
\[
 \frac1{p!q!}\int_0^1(1-s)^ps^q\,\dd s=\frac{1}{(p+q+1)!}.
\]
The beta integral is evaluated by induction on $p$: for $p=0$ it is
$1/(q+1)$, and integration by parts gives
$\int_0^1(1-s)^ps^q\dd s=\frac{p}{q+1}\int_0^1(1-s)^{p-1}s^{q+1}\dd s$, so
$\int_0^1(1-s)^ps^q\dd s=\frac{p!}{(q+1)(q+2)\cdots(q+p)}\cdot\frac1{q+p+1}
=\frac{p!q!}{(p+q+1)!}$. This proves \eqref{eq:duhamel}; interchanging sum
and integral is justified by the absolutely convergent majorant
$e^{(1-s)\norm A}\norm He^{s\norm A}$.

Multiplying \eqref{eq:duhamel} on the left by $e^{-A}$ gives
$\int_0^1e^{-sA}He^{sA}\dd s$, and \cref{thm:campbell} with $X=A$ and
parameter $-s$ turns the integrand into $e^{-s\ad_A}H$. Integrating the
absolutely convergent series $\sum_n(-s)^n\ad_A^nH/n!$ term by term gives
$\sum_n(-1)^n\ad_A^nH/(n+1)!$, which is \eqref{eq:leftdexp}. Multiplying
\eqref{eq:duhamel} on the right by $e^{-A}$ gives
$\int_0^1e^{(1-s)A}He^{-(1-s)A}\dd s=\int_0^1e^{u\ad_A}H\,\dd u$, and the same
integration gives \eqref{eq:rightdexp}.
\end{proof}
The quotient in $\phi(z)=(1-e^{-z})/z$ is notation for an entire function
with $\phi(0)=1$; \eqref{eq:leftdexp} does not require $\ad_A$ to be
invertible. The same statements hold intrinsically for finite-dimensional
Lie groups, where multiplication by $e^{-A}$ means left or right translation
of tangent vectors; an intrinsic derivation is given in \cref{sec:liegroups}.
In particular $(\dd\exp)_0=\id$.

\subsection{Bernoulli coefficients, including the convention at degree one}
The formal reciprocal of $\phi$ is
\begin{equation}\label{eq:betadef}
 \beta(z)=\frac{z}{1-e^{-z}}=\sum_{n\geq0}\bplus{n}z^n,\qquad
 \bplus{n}=\frac{B_n^+}{n!},
\end{equation}
which defines the Bernoulli numbers $B_n^+$ in the convention
$B_1^+=+\tfrac12$ used in the note of the source page. The other common
convention is generated by $\gamma(z)=z/(e^z-1)=\beta(-z)=\sum\bminus{n}z^n$,
with $\bminus{n}=(-1)^n\bplus{n}$ and $B_1^-=-\tfrac12$.

\begin{proposition}[Complete Bernoulli coefficients]\label{prop:bernoulli}
$\bplus{0}=1$, and for $n\geq1$ the coefficients are determined by the
recursion
\begin{equation}\label{eq:bernrec}
 \bplus{n}=-\sum_{j=1}^{n}\frac{(-1)^j}{(j+1)!}\,\bplus{n-j},
\end{equation}
and by the finite nonrecursive formula
\begin{equation}\label{eq:bernfinite}
 \bplus{n}=\sum_{k=1}^{n}(-1)^{n+k}
 \sum_{\substack{j_1+\cdots+j_k=n\\j_i\geq1}}\ \prod_{i=1}^{k}\frac1{(j_i+1)!}.
\end{equation}
Moreover
\begin{equation}\label{eq:coth}
 \beta(z)=\frac z2\Bigl(1+\coth\frac z2\Bigr)
 =1+\frac z2+\frac{z^2}{12}-\frac{z^4}{720}+\frac{z^6}{30240}
 -\frac{z^8}{1209600}+\cdots,
\end{equation}
so $\bplus{2m+1}=0$ for $m\geq1$; the even coefficients are
\begin{equation}\label{eq:bernzeta}
 \bplus{2m}=\frac{2(-1)^{m-1}\zeta(2m)}{(2\pi)^{2m}}\qquad(m\geq1),
\end{equation}
where $\zeta(2m)=\sum_{k\geq1}k^{-2m}$; and the Taylor series of $\beta$ has
radius of convergence exactly $2\pi$.
\end{proposition}
\begin{proof}
Write $\phi(z)=\sum_{j\geq0}\phi_jz^j$ with $\phi_j=(-1)^j/(j+1)!$, so
$\phi_0=1$. Comparing the coefficient of $z^n$, $n\geq1$, in
$\phi(z)\beta(z)=1$ gives $\sum_{j=0}^n\phi_j\bplus{n-j}=0$, which is
\eqref{eq:bernrec}. For \eqref{eq:bernfinite}, write $\phi=1+A$ with
$A(z)=\sum_{j\geq1}\phi_jz^j$ of positive order and expand
$\beta=(1+A)^{-1}=\sum_{k\geq0}(-A)^k$, a formally convergent geometric
series. The coefficient of $z^n$ in $(-A)^k$ is
$(-1)^k\sum_{j_1+\cdots+j_k=n}\prod_i\phi_{j_i}
=(-1)^k(-1)^n\sum\prod_i1/(j_i+1)!$, and summing over $k$ from $1$ to $n$
(no $k>n$ contributes because each $j_i\geq1$) gives \eqref{eq:bernfinite}.

Multiplying numerator and denominator of $\beta$ by $e^{z/2}$,
\[
 \beta(z)=\frac{ze^{z/2}}{e^{z/2}-e^{-z/2}}
 =\frac z2\cdot\frac{e^{z/2}+e^{-z/2}+(e^{z/2}-e^{-z/2})}{e^{z/2}-e^{-z/2}}
 =\frac z2\Bigl(\coth\frac z2+1\Bigr).
\]
The function $\frac z2\coth\frac z2=\beta(z)-\frac z2$ is even in $z$ and
regular at zero, so all its odd Taylor coefficients vanish; this proves
$\bplus{2m+1}=0$ for $m\geq1$. The displayed numerical coefficients follow
from \eqref{eq:bernrec}: $\bplus1=\tfrac12$, $\bplus2=\tfrac1{12}$,
$\bplus4=-\tfrac1{720}$, $\bplus6=\tfrac1{30240}$,
$\bplus8=-\tfrac1{1209600}$; for instance
$\bplus1=-\phi_1\bplus0=\tfrac12$ and
$\bplus2=-\bigl(\phi_1\bplus1+\phi_2\bplus0\bigr)
=-\bigl(-\tfrac12\cdot\tfrac12+\tfrac16\bigr)=\tfrac1{12}$.

For the even coefficients we derive the partial fraction expansion of
$\coth$ from an elementary Fourier series. Fix a real $z>0$ and let $g$ be
the $1$-periodic extension of $u\mapsto\cosh(z(u-\tfrac12))$ from $[0,1]$;
it is continuous, since $g(0)=g(1)=\cosh(z/2)$, and piecewise smooth. Its
Fourier coefficients are
\begin{align*}
 c_k&=\int_0^1\cosh\bigl(z(u-\tfrac12)\bigr)e^{-2\pi\ii ku}\,\dd u
 =\frac12e^{-z/2}\int_0^1e^{(z-2\pi\ii k)u}\dd u
  +\frac12e^{z/2}\int_0^1e^{-(z+2\pi\ii k)u}\dd u\\
 &=\frac{e^{z/2}-e^{-z/2}}{2(z-2\pi\ii k)}
  +\frac{e^{z/2}-e^{-z/2}}{2(z+2\pi\ii k)}
 =\frac{2z\sinh(z/2)}{z^2+4\pi^2k^2},
\end{align*}
using $e^{\pm2\pi\ii k}=1$. Since $c_k=O(k^{-2})$, the Fourier series
$\sum_{k\in\Z}c_ke^{2\pi\ii ku}$ converges absolutely and uniformly to a
continuous function with the same Fourier coefficients as $g$; by uniqueness
of Fourier coefficients for continuous functions, its sum is $g$. Evaluating
at $u=0$,
\[
 \cosh\frac z2=\frac{2\sinh(z/2)}{z}+\sum_{k\geq1}\frac{4z\sinh(z/2)}{z^2+4\pi^2k^2},
 \qquad\text{so}\qquad
 \coth\frac z2=\frac2z+4z\sum_{k\geq1}\frac1{z^2+4\pi^2k^2}.
\]
Both sides are meromorphic functions of $z\in\C$; the right side converges
normally on compact sets avoiding the points $\pm2\pi\ii k$, so the identity
proved for real $z>0$ extends to all $z$ away from the poles by the identity
theorem. Substituting into \eqref{eq:coth},
\begin{equation}\label{eq:partialfractions}
 \beta(z)=1+\frac z2+2z^2\sum_{k\geq1}\frac1{z^2+4\pi^2k^2}.
\end{equation}
For $|z|<2\pi$ expand each term geometrically,
$\frac{1}{z^2+4\pi^2k^2}=\sum_{m\geq0}\frac{(-1)^mz^{2m}}{(2\pi k)^{2m+2}}$;
the double series converges absolutely because
$\sum_k\sum_m|z|^{2m}/(2\pi k)^{2m+2}=\sum_k1/(4\pi^2k^2-|z|^2)<\infty$, so
the order of summation may be interchanged, giving
$\beta(z)=1+\frac z2+\sum_{m\geq0}\frac{2(-1)^m\zeta(2m+2)}{(2\pi)^{2m+2}}z^{2m+2}$,
which is \eqref{eq:bernzeta}. Finally, $1-e^{-z}$ vanishes exactly at
$2\pi\ii\Z$; at $z=0$ the zero of the denominator is cancelled by the
numerator, while at $\pm2\pi\ii$ the numerator is nonzero, so $\beta$ has
genuine poles there and nowhere closer to the origin. The Taylor radius is
therefore exactly $2\pi$. (Equivalently, \eqref{eq:bernzeta} gives
$|\bplus{2m}|^{1/2m}\to1/(2\pi)$ because $\zeta(2m)\to1$.)
\end{proof}
Functional notation such as $\beta(\ad_X)$ does not require $\ad_X$ to be
invertible. It means the formal series, or the analytic functional calculus in
a domain avoiding the poles; in a finite-dimensional space the latter is
defined whenever $\spec(\ad_X)$ avoids the nonzero points $2\pi\ii\ell$
(\cref{sec:regularX}).

\subsection{The function \texorpdfstring{$\psi$}{psi} and its complete expansion}
The other scalar function on the source page is
\begin{equation}\label{eq:psi}
 \psi(x)=\frac{x\log x}{x-1}=1-\sum_{n\geq1}\frac{(1-x)^n}{n(n+1)},
 \qquad|1-x|<1,\qquad \psi(1)=1,
\end{equation}
and, with the logarithm branch satisfying $\log(e^z)=z$ near $z=0$,
\begin{equation}\label{eq:psibeta}
 \psi(e^z)=\beta(z)=\sum_{n\geq0}B_n^+\frac{z^n}{n!}.
\end{equation}
\begin{proof}[Proof of \eqref{eq:psi} and \eqref{eq:psibeta}]
Put $u=1-x$ and use $\log(1-u)=-\sum_{j\geq1}u^j/j$ for $|u|<1$. Then
\[
 \psi(1-u)=\frac{(1-u)\log(1-u)}{-u}
 =(1-u)\sum_{j\geq1}\frac{u^{j-1}}{j}
 =\sum_{j\geq1}\frac{u^{j-1}}{j}-\sum_{j\geq1}\frac{u^{j}}{j}
 =1+\sum_{n\geq1}\Bigl(\frac1{n+1}-\frac1n\Bigr)u^n,
\]
and $\frac1{n+1}-\frac1n=-\frac1{n(n+1)}$. The same computation is valid as
an identity of formal power series in $u$. Substituting $x=e^z$ into
$x\log x/(x-1)$ gives $ze^z/(e^z-1)=z/(1-e^{-z})=\beta(z)$; as a formal
identity, substitute the positive-order series $u=1-e^z$ into
\eqref{eq:psi}.
\end{proof}
Equation \eqref{eq:psibeta} must not be read as a branch-independent
identity for arbitrary complex $z$; it is an identity of power series at
$z=0$, and analytically it holds wherever the logarithm branch used in
$\psi$ satisfies $\log(e^z)=z$.
